\section{Injecting a fault whose location is known}
\label{sec:method}

The missing ingredient is an answer. If we know which boundary is wrong, a trust
signal stops being something the system relies on and becomes something we can score.

We obtain that answer by injecting the fault ourselves. A fault is a small
transformation of one boundary function's source. It is applied by locating the
function, rewriting its abstract syntax tree, and splicing the result back through
the same scope validator a repair must pass. That last detail is what makes the
ground truth trustworthy rather than asserted: an injected fault is confined to one
boundary by the same mechanism that confines a repair, so if the confinement is
broken the injection fails rather than silently spreading.

Contamination is read from the pipeline rather than declared by us. A boundary is
downstream of the injected one when the injected function's result reaches its
arguments, either directly as a nested call or through a module-level name bound to
that result. Ground truth is therefore a pair: the injected boundary, and the
ordered list of boundaries its output reaches.

Two guards keep the corpus honest. An injection that leaves the source unchanged is
refused, because a fault nothing can detect would enter the corpus and depress every
signal's score equally, making every arm look worse without distinguishing them. And
a test asserts that every fault in the catalogue still changes something on a
representative pipeline; when we added it, it caught one that did not.

\subsection{The fault catalogue}

The catalogue separates faults observed in this system's own operational logs from
synthetic ones, because the two carry different weight. Three are observed. A field
disappears from a stage's output after upstream source reads failed, which this
system's logs record alongside HTTP 429, 403 and 500 responses from four source
registries. A collected sequence is shortened while the stage still reports success.
An identifier is replaced by a plausible placeholder, which is what a schema-bound
worker actually returned before a reviewing model discarded its output. Two are
synthetic: a comparison operator is inverted, and a ranking is reversed while every
value in it stays present and valid.

An invariant here is an acceptance criterion attached to a boundary and checked
deterministically against the value it produced; a concrete example is given in
Appendix~\ref{app:invariant}.

\subsection{What a trust signal is scored on}

A signal labels every boundary trusted or suspect. Scoring compares those labels
with the injected fault and reports four things.

\emph{False trust} is the injected boundary being labelled trusted. This is the
safety failure, and a signal that labels everything trusted scores perfectly on cost
and fails completely here.

\emph{False suspicion} is a boundary that is neither the injected one nor
contaminated by it being flagged. This is the cost, paid in redone work.

\emph{Localisation error} is signed, and the sign carries the finding. The resume
point is the earliest suspect boundary, which is where a run acting on these labels
would restart. A negative error means it restarts upstream of the fault and redoes
sound work. A positive error means it restarts past the fault and keeps output the
fault has already contaminated. Averaging the two would hide the second failure
inside the first, so we report them separately.

\emph{Shipped contaminated output} records whether the positive case occurred.

\subsection{Arms}

Four signals are in scope. Structural schema validation checks a stage's output
against its declared JSON schema. Content invariants check criteria the boundary
declared in advance: required fields, minimum counts, value patterns and ordering.
Both are deterministic and cost nothing per fault, which is what allows the corpus to
be large. The third arm is the reviewing model itself, which is the only arm that
costs tokens and the only one whose labels vary between runs on the same input.
